\chapter{Anemia}

\section*{Definition}
Anemia is a lower-than-expected hemoglobin concentration, reducing the blood's oxygen-carrying capacity. The threshold depends on age, sex, pregnancy, altitude, and other factors. It is a finding with many causes, and some people have few symptoms until anemia is severe.

\section*{When to See a Doctor}
\begin{itemize}
\item Unexplained fatigue, weakness, reduced exercise tolerance, or lethargy
\item Pallor (pale skin, conjunctiva, nail beds)
\item Shortness of breath on exertion, chest discomfort, fainting, or rapidly worsening symptoms
\item Dizziness or lightheadedness
\item Rapid heart rate or palpitations
\item Cold hands and feet
\item Brittle nails or spoon-shaped nails (koilonychia) suggesting iron deficiency
\item Glossitis (smooth, red tongue) suggesting B12 or folate deficiency
\item Pica (craving ice, dirt, or starch) suggesting iron deficiency
\end{itemize}

\section*{How It Develops}
Anemia results from reduced red-cell production, increased destruction (hemolysis), blood loss, or a combination. Reduced production can result from iron, vitamin B12, or folate deficiency; inflammation; kidney disease; endocrine illness; or bone-marrow disorders. Hemolysis may be due to inherited red-cell abnormalities or acquired immune, infectious, mechanical, or drug-related causes. Blood loss may be acute or chronic and can be hidden in the gastrointestinal or urinary tract.

\section*{Causes}
\begin{itemize}
\item Iron deficiency (most common cause worldwide)
\item Vitamin B12 deficiency (pernicious anemia, malabsorption)
\item Folate deficiency
\item Chronic disease/inflammation (anemia of chronic disease)
\item Acute or chronic blood loss (GI bleeding, heavy menstruation, trauma)
\item Hemolytic anemias (sickle cell disease, thalassemia, autoimmune hemolytic anemia)
\item Bone marrow disorders (aplastic anemia, myelodysplasia, leukemia)
\item Kidney disease (decreased erythropoietin production)
\item Medications and treatments, including chemotherapy, some antiretrovirals, and selected other drugs
\item Genetic disorders (G6PD deficiency, hereditary spherocytosis)
\end{itemize}

\section*{Related Symptoms}
\begin{itemize}
\item Fatigue
\item Shortness of breath
\item Pallor
\item Dizziness
\item Palpitations
\item Cold intolerance
\item Brittle nails
\item Headache
\item Restless legs syndrome (associated with iron deficiency)
\end{itemize}

\section*{Medical Evaluation}
\begin{itemize}
\item Complete blood count (CBC) with red cell indices
\item Serum ferritin and, when needed, iron, transferrin saturation, and related studies; ferritin can be falsely normal or high during inflammation
\item Serum vitamin B12 and folate levels
\item Reticulocyte count to assess bone marrow response
\item Peripheral blood smear for cell morphology
\item Hemoglobin electrophoresis if hemoglobinopathy is suspected
\item Evaluation for bleeding, including gastrointestinal evaluation in adults with iron-deficiency anemia when clinically indicated; bone-marrow examination is reserved for selected suspected marrow disorders
\end{itemize}

\section*{Treatment Options}
\begin{itemize}
\item Oral iron is commonly used for iron deficiency; intravenous iron is considered when oral iron is not tolerated, ineffective, not absorbed, or rapid replacement is needed
\item Vitamin B12 injections or high-dose oral B12 for B12 deficiency
\item Folate replacement for confirmed or strongly suspected folate deficiency, after considering vitamin B12 deficiency so that neurologic disease is not masked
\item Erythropoiesis-stimulating agents (ESA) for anemia of chronic kidney disease
\item Treatment of underlying chronic condition (infection, inflammation, malignancy)
\item Blood transfusion for selected patients with severe or symptomatic anemia, active bleeding, or inadequate oxygen delivery; the decision depends on the hemoglobin, symptoms, rate of fall, and overall condition
\item Specialist treatment for hemolytic or inherited anemias; splenectomy is used only in selected cases after weighing long-term risks
\end{itemize}

\section*{Self-Care and Home Management}
\begin{itemize}
\item Consume iron-rich foods (red meat, spinach, lentils, fortified cereals)
\item Pair iron-rich foods with vitamin C (citrus, bell peppers) to enhance absorption
\item If taking oral iron, ask about interactions with tea, coffee, calcium, antacids, and other medicines; dietary advice should fit the person's overall nutrition
\item Eat foods rich in B12 (meat, fish, eggs, dairy) and folate (leafy greens, beans)
\item Rest when fatigued and seek urgent care for chest pain, fainting, severe breathlessness, black or bloody stools, or heavy bleeding
\item Avoid over-the-counter iron supplements without medical guidance
\item Limit alcohol consumption, which can interfere with folate metabolism
\end{itemize}

\section*{References}
\begin{thebibliography}{99}
\bibitem{anemia:ref1} Camaschella C. ``Iron-deficiency anemia.'' \textit{N Engl J Med}. 2015;372(19):1832-1843.
\bibitem{anemia:ref2} Green R, Datta Mitra A, et al. ``Vitamin B12 deficiency.'' \textit{Nat Rev Dis Primers}. 2017;3:17040.
\bibitem{anemia:ref3} Short MW, Domagalski JE. ``Iron deficiency anemia: evaluation and management.'' \textit{Am Fam Physician}. 2013;87(2):98-104.
\bibitem{anemia:ref4} Kaushansky K, Lichtman MA, et al. \textit{Williams Hematology}, 10th ed. McGraw-Hill, 2021.
\bibitem{anemia:ref5} UpToDate. ``Diagnostic approach to anemia in adults.'' Wolters Kluwer, 2023.
\bibitem{anemia:ref6} World Health Organization. Guideline on haemoglobin cutoffs to define anaemia in individuals and populations. 2024.
\bibitem{anemia:ref7} Snook J, Bhala N, Beales ILP, et al. British Society of Gastroenterology guidelines for the management of iron deficiency anaemia in adults. \textit{Gut}. 2021;70:2030--2051.
\end{thebibliography}
